\chapter{Chapter 2 --- Background and Related
Work}\label{chapter-2-background-and-related-work}

\begin{quote}
\textbf{Citation status.} Every reference in this chapter was checked
against its publisher record or the source itself during preparation;
\S{}2.9 records the outcome. \textbf{One cited work could not be located
and has been removed} --- see \S{}2.9. Page-level details (exact page
numbers, section numbers) should still be confirmed against the PDFs
when the bibliography is typeset.
\end{quote}

\section{2.1 Just-in-time software defect
prediction}\label{just-in-time-software-defect-prediction}

\subsection{2.1.1 Task definition}\label{task-definition}

JIT-SDP predicts, for each incoming commit, whether it introduces a
defect. Formally, given a commit \emph{c} with feature vector \emph{x} \ensuremath{\in}
\ensuremath{\mathbb{R}}\^{}d available at commit time, predict \emph{y} \ensuremath{\in} \{0, 1\}.

The task's appeal is timing. A prediction available at commit time
reaches the developer who wrote the change while the change is still
current, which is the difference between a useful signal and a
retrospective one.

\subsection{2.1.2 The Kamei feature set}\label{the-kamei-feature-set}

Kamei et al.~established the modern formulation, together with the
change-level metrics that remain standard. The 14 features group into
five families:

\begin{longtable}[]{@{}
  >{\raggedright\arraybackslash}p{(\linewidth - 4\tabcolsep) * \real{0.3333}}
  >{\raggedright\arraybackslash}p{(\linewidth - 4\tabcolsep) * \real{0.3333}}
  >{\raggedright\arraybackslash}p{(\linewidth - 4\tabcolsep) * \real{0.3333}}@{}}
\toprule\noalign{}
\begin{minipage}[b]{\linewidth}\raggedright
Family
\end{minipage} & \begin{minipage}[b]{\linewidth}\raggedright
Features
\end{minipage} & \begin{minipage}[b]{\linewidth}\raggedright
Intuition
\end{minipage} \\
\midrule\noalign{}
\endhead
\bottomrule\noalign{}
\endlastfoot
\textbf{Diffusion} & NS, ND, NF, Entropy & A change spread across many
subsystems, directories and files is harder to reason about \\
\textbf{Size} & LA, LD, LT & Large changes carry more opportunity for
error \\
\textbf{Purpose} & FIX & Changes that fix defects are themselves
defect-prone \\
\textbf{History} & NDEV, AGE, NUC & Code touched by many developers,
recently, and often \\
\textbf{Experience} & EXP, REXP, SEXP & Author familiarity with the code
being changed \\
\end{longtable}

Kamei et al.~report approximately \textbf{68\% accuracy and 64\% recall}
on defect-inducing changes. These features are used throughout this
thesis, taken as published rather than recomputed so that feature
extraction cannot differ between the reference labels and the SZZ
variants.

\subsection{2.1.3 Class imbalance and concept
drift}\label{class-imbalance-and-concept-drift}

Defect-introducing commits are a small minority --- \textbf{8.54\%} on
the corpus studied here. Two consequences run through this thesis.

\textbf{Accuracy is uninformative.} A model predicting ``clean'' for
every commit scores 91.5\% accuracy on this corpus. \S{}2.6 covers metric
selection.

\textbf{Oversampling is standard, and it is not free.} Rebalancing
methods increase the influence of minority examples, which is exactly
what is wanted when those examples are correct. Chapter 6 shows what
happens when they are not.

Projects also change over time --- team composition, architecture,
process --- so the relationship between features and defect-proneness is
not stationary. Concept drift in this setting is addressed by Cabral and
Minku and by subsequent work; \S{}2.4 returns to it.

\section{2.2 SZZ and its variants}\label{szz-and-its-variants}

\subsection{2.2.1 The original algorithm}\label{the-original-algorithm}

\'{S}liwerski, Zimmermann and Zeller, in \emph{When Do Changes Induce
Fixes?} (MSR 2005), introduced the procedure now universally called SZZ
after their initials. Given a commit identified as fixing a defect, it:

\begin{enumerate}
\def\labelenumi{\arabic{enumi}.}
\tightlist
\item
  Extracts the lines modified by the fix;
\item
  Applies \texttt{git\ blame} (or equivalent) to find, for each line,
  the most recent prior commit to have modified it;
\item
  Labels those prior commits defect-introducing.
\end{enumerate}

The algorithm is an inference, and it rests on two premises: that the
lines a fix modifies are the lines that were defective, and that the
last commit to touch a line is the commit that made it wrong.

\subsection{2.2.2 The failure taxonomy}\label{the-failure-taxonomy}

\textbf{Tangled commits} violate the first premise. Bug-fixing commits
routinely contain reformatting, renaming, comment updates and test
changes alongside the fix. Herbold et al., using four annotators per
line, found that \textbf{between 17\% and 32\% of all changes in
bug-fixing commits address the underlying problem}, rising to
\textbf{66--87\%} when only production-code files are considered; about
\textbf{11\%} of lines produced active annotator disagreement. Every
tangled line is blamed, and every commit that last touched one is
mislabelled.

\textbf{Blame attribution error} violates the second premise.
\texttt{git\ blame} identifies the last commit to \emph{touch} a line,
which need not be the commit that made it \emph{wrong}. A later cosmetic
edit takes the blame from the change that introduced the fault. This
thesis refers to this as the \textbf{syntactic line-blame fallacy}, and
\S{}2.2.4 explains why it is not measurable in this study's design.

\textbf{Structural invisibility.} A defect introduced by omission ---
code that was never written --- has no line for blame to find. This
mechanism is real in general but, as Chapter 4 explains, is \emph{not
available} as an explanation within this thesis's corpus.

\subsection{2.2.3 The variants}\label{the-variants}

Each variant adds filtering or selection intended to remove B-SZZ's
over-collection:

\begin{longtable}[]{@{}
  >{\raggedright\arraybackslash}p{(\linewidth - 4\tabcolsep) * \real{0.3333}}
  >{\raggedright\arraybackslash}p{(\linewidth - 4\tabcolsep) * \real{0.3333}}
  >{\raggedright\arraybackslash}p{(\linewidth - 4\tabcolsep) * \real{0.3333}}@{}}
\toprule\noalign{}
\begin{minipage}[b]{\linewidth}\raggedright
Variant
\end{minipage} & \begin{minipage}[b]{\linewidth}\raggedright
Mechanism
\end{minipage} & \begin{minipage}[b]{\linewidth}\raggedright
Attribution
\end{minipage} \\
\midrule\noalign{}
\endhead
\bottomrule\noalign{}
\endlastfoot
\textbf{B-SZZ} & The original & \'{S}liwerski et al., MSR 2005 \\
\textbf{AG-SZZ} & Annotation graphs for precise line tracing; filters
cosmetic changes such as blank lines and comments & Kim et al., 2006 \\
\textbf{MA-SZZ} & Additionally excludes meta-changes --- branch merges,
property changes --- that do not alter program behaviour & da Costa et
al., 2017 \\
\textbf{R-SZZ} & Retains only the most recent candidate commit & Davies
et al., 2014 \\
\textbf{L-SZZ} & Retains only the candidate with the most changed lines
& Davies et al., 2014 \\
\textbf{RA-SZZ} & Refactoring-aware; detects and filters refactorings
via RefDiff and RefactoringMiner (Java only) & Neto et al., 2018 \\
\end{longtable}

The progression assumes that B-SZZ's dominant error is over-collection
and that filtering it out improves the labels. \textbf{Chapter 4 tests
that assumption directly and finds that for the most aggressive
variants, more than half of what the filters remove is correct.}

\subsection{2.2.4 Developer-informed
oracles}\label{developer-informed-oracles}

Evaluating SZZ requires a reference that does not itself use SZZ. Rosa
et al.~constructed one by mining commit messages in which developers
\emph{explicitly name} the commit that introduced the defect being fixed
--- yielding \textbf{1,930 fix-to-inducing links across eight
programming languages}, extended in the journal version to an evaluation
of nine variants over \textbf{2,304 instances}. Against that oracle,
\textbf{R-SZZ performed best (F-measure \ensuremath{\approx} 61\%, precision \ensuremath{\approx} 66\%)} while
B-SZZ achieved higher recall (\ensuremath{\approx} 0.69) at markedly lower precision (\ensuremath{\approx}
0.39).

\textbf{Two points about this dataset matter for the present work.}

First, it \textbf{cannot train a JIT-SDP model.} It contains
fix-to-inducing links and no clean commits, so there is no negative
class and no complete labelling of project history.

Second, \textbf{its precision figures are not comparable to this
thesis's without reconciling the denominator.} Rosa et al.~score
\emph{conditionally} on bug-fixing commits already known to have an
inducing commit. This thesis scores over an entire project history of
27,319 commits, the overwhelming majority of which are clean. A variant
that is right two times in three when asked ``which earlier commit
introduced this known defect?'' can still be wrong four times in five
when asked ``is this commit, drawn from all of history,
defect-introducing?'' \textbf{The difference is large enough to read as
a contradiction, and Chapter 8 addresses it explicitly.}

\subsection{2.2.5 JIT-Defects4J and the reference
labels}\label{jit-defects4j-and-the-reference-labels}

Ni et al.~published \textbf{JIT-Defects4J}, an extension of LLTC4J,
comprising \textbf{27,319 commits from 21 Java projects, of which 2,332
(8.54\%) are labelled defect-introducing}. Its construction is
two-stage: lines within bug-fixing commits were annotated manually,
retained where at least three participants agreed; \texttt{git\ blame}
was then applied to those verified lines.

\textbf{The labels are therefore \texttt{git\ blame} seeded with
human-verified fix lines --- not manually verified ground truth for
defect-introducing commits.} This thesis uses them as its reference
precisely because the distinction is tractable: they control for
tangling, the dominant false-positive mechanism, while sharing SZZ's
blame step. Comparisons against them isolate the cost of tangled
commits, and blame error, being common to both sides, cancels. Chapter 3
\S{}3.2.2 develops the consequences.

\subsection{2.2.6 Other JIT-SDP datasets}\label{other-jit-sdp-datasets}

\begin{longtable}[]{@{}
  >{\raggedright\arraybackslash}p{(\linewidth - 6\tabcolsep) * \real{0.2500}}
  >{\raggedright\arraybackslash}p{(\linewidth - 6\tabcolsep) * \real{0.2500}}
  >{\raggedright\arraybackslash}p{(\linewidth - 6\tabcolsep) * \real{0.2500}}
  >{\raggedright\arraybackslash}p{(\linewidth - 6\tabcolsep) * \real{0.2500}}@{}}
\toprule\noalign{}
\begin{minipage}[b]{\linewidth}\raggedright
Dataset
\end{minipage} & \begin{minipage}[b]{\linewidth}\raggedright
Scale
\end{minipage} & \begin{minipage}[b]{\linewidth}\raggedright
Labels
\end{minipage} & \begin{minipage}[b]{\linewidth}\raggedright
Relevance here
\end{minipage} \\
\midrule\noalign{}
\endhead
\bottomrule\noalign{}
\endlastfoot
\textbf{ApacheJIT} & 106,674 commits (28,239 bug-inducing), 14 Apache
projects & SZZ-derived & Larger, but cannot serve as an oracle for a
study of SZZ noise \\
\textbf{Defectors} & \textasciitilde213K files, 24 Python projects, 18
domains & SZZ-derived & File-level; same objection \\
\textbf{ReDef} & 3,164 defective / 10,268 clean functions, 22 C/C++
projects & \textbf{Revert-anchored}, LLM-triaged, \textasciitilde92\%
audited precision & Non-blame labels, but function-level, C/C++, no
Kamei features, no fix dates \\
\end{longtable}

\textbf{Scale is not the binding constraint for this thesis; label
provenance is.} A study measuring what SZZ noise costs cannot use SZZ
labels as its reference, which excludes the two largest datasets
available.

\section{2.3 Evaluation methodology in defect
prediction}\label{evaluation-methodology-in-defect-prediction}

\subsection{2.3.1 Temporal leakage}\label{temporal-leakage}

Random \emph{k}-fold cross-validation on temporally ordered data allows
a model to train on future commits and be tested on past ones.
Time-aware evaluation --- training only on the past and evaluating only
on the future --- has been advocated repeatedly, notably by Falessi and
colleagues, who caution researchers against ``time travel'' and against
generalising beyond the evaluation context. The reported effect is
consistent in direction with this thesis's: in-sample performance is
systematically inflated relative to time-aware performance. Random
\emph{k}-fold nonetheless remains widespread in the JIT-SDP literature.

\textbf{Chapter 5 measures the consequence rather than restating the
argument.} The inflation is +0.137 MCC on a high-capacity model and it
scales with model capacity, which means comparisons \emph{between} model
families under this protocol are confounded in a direction that favours
larger models.

\subsection{2.3.2 Prequential evaluation}\label{prequential-evaluation}

For streams, the standard protocol is \emph{test-then-train}: each
example is first predicted, then used for learning. Summarising a metric
over a stream requires a forgetting mechanism, conventionally a fading
factor applied to the error estimate. Gama, Sebasti\~{a}o and Rodrigues
(\emph{On Evaluating Stream Learning Algorithms}, Machine Learning,
2013) defend prequential error with forgetting and prove that,
\textbf{for consistent learning algorithms on stationary data, the
prequential error computed with fading factors converges to the Bayes
error}.

\textbf{An important caveat for this thesis.} That convergence result is
established for an \emph{error rate} --- a decomposable loss accumulated
example by example. \textbf{MCC is not decomposable}, being a non-linear
function of all four confusion-matrix cells, so neither the construction
nor its convergence guarantee extends to it directly. This thesis
therefore reports two ad-hoc summaries of the MCC trajectory --- a
terminal fading value and a trajectory mean --- treats the trajectory
mean as primary \textbf{on measured variance rather than by appeal to a
standard}, and states explicitly that neither is canonical. An earlier
version of this work described one of them as ``the standard
estimator''; that was incorrect and is corrected in Chapter 3 \S{}3.4.2.

\subsection{2.3.3 Verification latency}\label{verification-latency}

Cabral and Minku identified verification latency as a first-class
concern in JIT-SDP: labels arrive with a delay, and the delay differs
systematically between classes, since a clean commit is only ever
\emph{presumed} clean while a defective one must be discovered. Their
framing of latency and of the resulting \textbf{class imbalance
evolution} is the setting this thesis adopts.

On the corpus studied here the median latency is \textbf{113 days}, with
\textbf{53\% of defect labels arriving after a 90-day window}. Under a
fixed window, a defect label that arrives late is delivered first as a
\emph{wrong clean label} and corrected later --- so latency is not
merely delay but active misinformation, and the misinformation concerns
exactly the minority class.

\section{2.4 Online learning for
JIT-SDP}\label{online-learning-for-jit-sdp}

\textbf{Online Bagging} (Oza and Russell, AISTATS 2001) adapts bagging
to streams by presenting each arriving example to each ensemble member
\emph{k} \textasciitilde{} Poisson(\ensuremath{\lambda}) times, achieving performance
comparable to batch bagging without storing the dataset.

\textbf{Oversampling Online Bagging (OOB)} (Wang, Minku and Yao) handles
imbalance by setting \ensuremath{\lambda} from the running observed class ratio --- tracked
with time-decayed metrics --- so minority examples are presented more
often as the imbalance grows. OOB and its undersampling counterpart UOB
are the standard data-level approaches for class-imbalanced streams.

\textbf{Oversampling Rate Boosting (ORB)} extends OOB by monitoring the
model's own recent predictions: when predictions are considerably biased
toward one class, the resampling rate of the opposite class is boosted.
Cabral and Minku also adopt a safety mechanism intended to avoid
training on potentially noisy or outlier minority-class examples.

\begin{quote}
\textbf{This machinery is the subject of Chapter 6's central finding.}
ORB's \ensuremath{\lambda} is a function of the \emph{observed rate} of defect labels ---
that is, of label \textbf{quantity}. It has no access to label
\textbf{quality}. Chapter 6 measures the coupling directly (Spearman \ensuremath{\rho} =
\ensuremath{-}1.000 between delivered defect-label supply and applied \ensuremath{\lambda}) and shows
that the mechanism correcting class imbalance is the mechanism that
amplifies false positives, hardest exactly when defect labels are
scarcest.
\end{quote}

Cabral and Minku's later work (\emph{Towards Reliable Online
Just-in-Time Software Defect Prediction}) analyses concept drift in
JIT-SDP across p(y), p(x\textbar y) and p(x), reports \textbf{defect
discovery delays ranging from one day to more than eleven years}, and
proposes a drift-adaptive method that monitors predictions on unlabelled
data rather than waiting for labels. \textbf{Their finding that medians
are typically at or below 90 days, and that a 90-day waiting period
trades off correct labelling against drift, is the basis for this
thesis's choice of \emph{W} = 90 days.} \textbf{The distinction from the
present work is the origin of the noise.} Latency-induced noise is a
timing artefact: the label is correct but late. SZZ-induced noise is a
\emph{labelling} artefact: the label is wrong and will never be
corrected. This thesis addresses the second, under conditions that
include the first.

\section{2.5 The weak-baseline problem}\label{the-weak-baseline-problem}

Deep learning approaches to JIT-SDP --- \textbf{DeepJIT} (Hoang et al.,
MSR 2019), which encodes commits hierarchically from line to hunk to
file to commit using a convolutional text model, and \textbf{CC2Vec}
(Hoang et al., ICSE 2020), which learns code-change representations
supervised by commit logs via a hierarchical attention network ---
reported substantial gains over feature-based models. Two subsequent
results complicate that picture.

Zeng et al., evaluating on over \textbf{310,370 changes}, found that
CC2Vec could not consistently outperform DeepJIT and that neither
consistently outperformed traditional feature-based JIT prediction. They
then constructed \textbf{LApredict} --- logistic regression on the
added-lines feature alone --- which outperformed both deep models while
being roughly \textbf{81,000--120,000\ensuremath{\times} faster} to train and test.

Pornprasit and Tantithamthavorn reported \textbf{JITLine}, a simpler
feature-based approach that was \textbf{26--38\% better in F-measure,
17--51\% more cost-effective, and 70--100\ensuremath{\times} faster} than CC2Vec and
DeepJIT.

\textbf{Both models are used in this thesis, and the reason is
methodological rather than deferential.} LApredict has minimal capacity;
JITLine has a great deal. Chapter 5 shows that the inflation
attributable to evaluation protocol is \textbf{+0.137 MCC for JITLine
and +0.034 for LApredict}, and that the self-scoring gap is significant
for all six SZZ variants under JITLine but only one under LApredict.
\textbf{Capacity is not incidental to the measurement problem --- it is
the variable that determines its size}, which offers an alternative
reading of why simple baselines have proven so hard to beat.

\section{2.6 Metric selection}\label{metric-selection}

Under an 8.54\% positive rate, accuracy is uninformative. This thesis
uses:

\textbf{Matthews Correlation Coefficient (MCC)} --- introduced by
Matthews (1975) for protein secondary-structure prediction and since
adopted widely for imbalanced binary classification --- as primary. It
incorporates all four confusion-matrix cells and is near zero for any
uninformed strategy. A majority-class predictor scores 0.0 MCC and
91.5\% accuracy on this corpus.

\textbf{G-mean}, the geometric mean of per-class recalls, which
collapses toward zero if either class is abandoned.

\textbf{Cohen's \ensuremath{\kappa}} (Cohen, 1960), the proportion of agreement corrected
for chance, for agreement between labelling procedures, where the
question is concordance rather than prediction quality.

Effort-aware metrics --- recall at a fixed inspection budget ---
arguably match the practitioner's objective more closely and are
identified as future work in Chapter 9, not reported here.

\section{2.7 Learning with noisy
labels}\label{learning-with-noisy-labels}

\textbf{Loss correction.} Natarajan et al.~(\emph{Learning with Noisy
Labels}, NIPS 2013) study binary classification under random
classification noise whose \textbf{flip probability is
class-conditional} --- precisely the noise model Chapter 4 measures ---
and construct a simple \textbf{unbiased estimator of any loss}, with
performance bounds for empirical risk minimisation on noisily labelled
i.i.d. data. The correction becomes numerically unstable as \ensuremath{\rho}\ensuremath{_0} + \ensuremath{\rho}\ensuremath{_1} \ensuremath{\rightarrow} 1.
\textbf{This matters directly here}: Chapter 4 measures \ensuremath{\rho}\ensuremath{_0} + \ensuremath{\rho}\ensuremath{_1}
approaching 0.8 for the conservative variants, which is why the
loss-correction arm evaluated in this thesis is capped.

\textbf{Sample selection.} Co-teaching (Han et al., NeurIPS 2018) trains
two networks simultaneously; on each mini-batch, each selects the
examples it considers clean and back-propagates only the examples
selected by its peer, on the premise that mislabelled examples exhibit
higher loss early in training. It assumes a batch setting with multiple
epochs and two models.

\textbf{Confident Learning} (Northcutt, Jiang and Chuang, JAIR 2021)
assumes a \textbf{class-conditional} noise process and directly
estimates the joint distribution between noisy and true labels, then
prunes, counts with probabilistic thresholds, and ranks examples by
confidence. It is a batch method requiring cross-validated out-of-sample
predictions over a fixed dataset.

\textbf{Why these do not transfer directly.} Each assumes a fixed
dataset, repeated passes, and labels that are wrong but \emph{stable}. A
JIT-SDP stream offers none of these: a single pass, labels arriving
late, and labels that are \textbf{wrong and then corrected}. Chapter 7
adapts the confidence-based idea to the streaming setting, and its
central negative result is informative about the whole family: a
confidence signal cannot distinguish \emph{``this label is wrong''} from
\emph{``my model has not learned this pattern yet''}, and the second
case is common under severe imbalance with a still-learning model.

\section{2.8 Positioning}\label{positioning}

\begin{longtable}[]{@{}
  >{\raggedright\arraybackslash}p{(\linewidth - 10\tabcolsep) * \real{0.1667}}
  >{\raggedright\arraybackslash}p{(\linewidth - 10\tabcolsep) * \real{0.1667}}
  >{\raggedright\arraybackslash}p{(\linewidth - 10\tabcolsep) * \real{0.1667}}
  >{\raggedright\arraybackslash}p{(\linewidth - 10\tabcolsep) * \real{0.1667}}
  >{\raggedright\arraybackslash}p{(\linewidth - 10\tabcolsep) * \real{0.1667}}
  >{\raggedright\arraybackslash}p{(\linewidth - 10\tabcolsep) * \real{0.1667}}@{}}
\toprule\noalign{}
\begin{minipage}[b]{\linewidth}\raggedright
Work
\end{minipage} & \begin{minipage}[b]{\linewidth}\raggedright
Imbalance
\end{minipage} & \begin{minipage}[b]{\linewidth}\raggedright
Drift
\end{minipage} & \begin{minipage}[b]{\linewidth}\raggedright
Latency
\end{minipage} & \begin{minipage}[b]{\linewidth}\raggedright
SZZ noise \textbf{measured}
\end{minipage} & \begin{minipage}[b]{\linewidth}\raggedright
Online
\end{minipage} \\
\midrule\noalign{}
\endhead
\bottomrule\noalign{}
\endlastfoot
Kamei et al. & \ensuremath{\circ} & \ensuremath{\circ} & \ensuremath{\circ} & \ensuremath{\circ} & \ensuremath{\circ} \\
Cabral \& Minku (ORB; drift work) & \ensuremath{\bullet} & \ensuremath{\bullet} & \ensuremath{\bullet} & \ensuremath{\circ} & \ensuremath{\bullet} \\
Rosa et al.~(developer-informed oracle) & \ensuremath{\circ} & \ensuremath{\circ} & \ensuremath{\circ} & \ensuremath{\bullet} \emph{(labels
only)} & \ensuremath{\circ} \\
Herbold et al.~(tangling) & \ensuremath{\circ} & \ensuremath{\circ} & \ensuremath{\circ} & \ensuremath{\bullet} \emph{(labels only)} & \ensuremath{\circ} \\
Zeng et al.~/ Pornprasit \& Tantithamthavorn & \ensuremath{\bullet} & \ensuremath{\circ} & \ensuremath{\circ} & \ensuremath{\circ} & \ensuremath{\circ} \\
Natarajan / Co-teaching / Confident Learning & \ensuremath{\bullet} & \ensuremath{\circ} & \ensuremath{\circ} & \ensuremath{\odot}
\emph{(noise model assumed, not measured)} & \ensuremath{\circ} \\
\textbf{This thesis} & \ensuremath{\bullet} & \ensuremath{\bullet} & \ensuremath{\bullet} & \ensuremath{\bullet} \textbf{downstream cost} & \ensuremath{\bullet} \\
\end{longtable}

\emph{(\ensuremath{\bullet} addressed \ensuremath{\cdot} \ensuremath{\odot} partially \ensuremath{\cdot} \ensuremath{\circ} not addressed)}

\textbf{The distinguishing cell is the fourth.} SZZ evaluation work
measures the noise and stops at the labels. Online JIT-SDP work operates
under latency but treats label noise as a nuisance to dampen rather than
a quantity to measure. Noisy-label learning assumes a noise model rather
than estimating one, and assumes a batch setting. This thesis measures
\textbf{what SZZ-origin noise costs downstream, under latency-aware
online evaluation}, and then tests whether that cost can be recovered.

The methodological device that makes this possible is small and worth
naming: \textbf{an independently constructed reference label set, held
out from training and used only for scoring.} Without it, a model
trained on SZZ can only be evaluated against SZZ, and the largest effect
in this thesis --- the +0.230 MCC self-scoring gap --- is invisible by
construction.

\section{2.9 Citation verification
status}\label{citation-verification-status}

Every work cited in this chapter was checked against its publisher
record or the source during preparation. The table records what was
confirmed.

\begin{longtable}[]{@{}
  >{\raggedright\arraybackslash}p{(\linewidth - 4\tabcolsep) * \real{0.3333}}
  >{\raggedright\arraybackslash}p{(\linewidth - 4\tabcolsep) * \real{0.3333}}
  >{\raggedright\arraybackslash}p{(\linewidth - 4\tabcolsep) * \real{0.3333}}@{}}
\toprule\noalign{}
\begin{minipage}[b]{\linewidth}\raggedright
Work
\end{minipage} & \begin{minipage}[b]{\linewidth}\raggedright
Venue
\end{minipage} & \begin{minipage}[b]{\linewidth}\raggedright
Confirmed
\end{minipage} \\
\midrule\noalign{}
\endhead
\bottomrule\noalign{}
\endlastfoot
\'{S}liwerski, Zimmermann \& Zeller, \emph{When Do Changes Induce Fixes?} &
MSR 2005 & Original SZZ; fix-inducing change analysis over CVS
archives \\
Kim et al. & 2006 & AG-SZZ: annotation graphs, cosmetic-change
filtering \\
da Costa et al. & 2017 & MA-SZZ: excludes meta-changes (branch merges,
property changes) \\
Davies et al. & 2014 & R-SZZ (most recent candidate), L-SZZ (most
changed lines) \\
Neto et al. & 2018 & RA-SZZ: refactoring-aware via RefDiff /
RefactoringMiner, Java only \\
Rosa et al. & ICSE 2021; \emph{JSS} 202:111729, 2023 & 1,930 links, 8
languages, 2,304 instances; R-SZZ best (F \ensuremath{\approx} 61\%, P \ensuremath{\approx} 66\%); B-SZZ
recall \ensuremath{\approx} 0.69, precision \ensuremath{\approx} 0.39; PySZZ \\
Herbold et al. & \emph{EMSE}, 2022 & 17--32\% of changes / 66--87\%
production code; four annotators per line; \textasciitilde11\%
disagreement \\
Ni et al. & \emph{ESEC/FSE} 2022 & JIT-Defects4J; 27,319 commits / 2,332
defective / 21 projects; two-stage construction \\
Kamei et al. & \emph{IEEE TSE} 39(6):757--773, 2013 & Task formulation,
change-level metrics; \textasciitilde68\% accuracy, 64\% recall \\
Cabral \& Minku & \emph{ICSE} 2019 & Class imbalance evolution,
verification latency, ORB, safety mechanism \\
Cabral \& Minku, \emph{Towards Reliable Online JIT-SDP} & --- & Drift
across p(y), p(x\textbar y), p(x); delays from 1 day to \textgreater11
years; medians at or below 90 days; 90-day waiting time as the
trade-off \\
Oza \& Russell, \emph{Online Bagging and Boosting} & AISTATS 2001, PMLR
R3:229--236 & Poisson(\ensuremath{\lambda}) online bagging matching batch performance \\
Wang, Minku \& Yao & --- & OOB / UOB; time-decayed class-imbalance
tracking \\
Gama, Sebasti\~{a}o \& Rodrigues, \emph{On Evaluating Stream Learning
Algorithms} & \emph{Machine Learning}, 2013 & Prequential error with
fading factors converges to Bayes error for consistent learners on
stationary data \\
Hoang et al., \emph{DeepJIT} & MSR 2019 & Hierarchical
line\ensuremath{\rightarrow}hunk\ensuremath{\rightarrow}file\ensuremath{\rightarrow}commit encoding, textCNN \\
Hoang et al., \emph{CC2Vec} & ICSE 2020 & Code-change representations
supervised by commit logs, hierarchical attention \\
Zeng et al. & \emph{ISSTA} 2021 & 310,370 changes; LApredict; 81k--120k\ensuremath{\times}
speedup \\
Pornprasit \& Tantithamthavorn, \emph{JITLine} & \emph{MSR} 2021,
pp.~369--379 & 26--38\% F-measure, 17--51\% cost-effectiveness, 70--100\ensuremath{\times}
faster \\
Keshavarz \& Nagappan, \emph{ApacheJIT} & \emph{MSR} 2022 & 106,674
commits, 28,239 bug-inducing, 14 projects \\
Mahbub, Shuvo \& Rahman, \emph{Defectors} & \emph{MSR} 2023 &
\textasciitilde213K files, 24 Python projects, 18 domains \\
ReDef & preprint, 2025 & 3,164 / 10,268 function-level changes, 22 C/C++
projects, revert-anchored, \textasciitilde92\% precision \\
Matthews & \emph{BBA} 405(2):442--451, 1975 & MCC, originally for
protein secondary-structure prediction \\
Cohen & \emph{Educ. Psychol. Meas.} 20:37--46, 1960 & \ensuremath{\kappa} as agreement
corrected for chance \\
Natarajan et al., \emph{Learning with Noisy Labels} & NIPS 2013 &
Class-conditional flip probability; unbiased loss estimator; ERM
bounds \\
Han et al., \emph{Co-teaching} & NeurIPS 2018 & Two networks exchanging
small-loss examples per mini-batch \\
Northcutt, Jiang \& Chuang, \emph{Confident Learning} & \emph{JAIR} 70,
2021 & Class-conditional noise; joint distribution estimation; prune,
count, rank \\
Zhao, Damevski \& Chen & \emph{ACM CSUR} 55(10):1--35, 2023 & Systematic
survey of JIT-SDP; 67 studies \\
Destefanis et al., \emph{An Audit of ML Experiments on Software Defect
Prediction} & \emph{EMSE}, 2026 & 101 audited papers; SCOPUS 2019--2023;
design, analysis and reporting practice \\
Falessi et al. & --- & Time-aware evaluation; caution against ``time
travel''; in-sample performance systematically inflated \\
\end{longtable}

\subsection{Two notes on the
bibliography}\label{two-notes-on-the-bibliography}

\textbf{One cited work was removed.} The project's planning documents
referenced ``Song et al.~(2022)'' as a method addressing latency-induced
label noise in JIT-SDP. \textbf{That work could not be located.}
Searching the area returns Cabral and Minku on verification latency and
Tabassum, Minku and colleagues on waiting-time labelling, either of
which may have been what was intended. The reference is therefore
omitted rather than guessed at, and \S{}2.4's claim about latency-noise
methods is attributed only to work that was confirmed. \textbf{If a
specific Song et al.~paper was intended, locate it before reinstating
the citation.}

\textbf{One attribution needs confirming.} A deep-learning-in-SDP
systematic review and meta-analysis exists (\emph{Information and
Software Technology}, 2023) and is the work the planning documents meant
by ``Zain et al.'', but the first-author attribution was not confirmed.
It is not cited in this chapter; if it is added, check the author list.

\textbf{One correction to the project's own outline.} An earlier outline
attributed to Herbold et al.~a finding that ``\textasciitilde50\%'' of
SZZ labels are wrong. The verified figures are \textbf{17--32\% of all
changes in bug-fixing commits} addressing the underlying problem, rising
to \textbf{66--87\%} for production-code files only --- a different
quantity, measured at line rather than commit level. The
``\textasciitilde50\%'' figure does not appear in the source and must
not be cited.

\textbf{Remaining work for the bibliography.} Page numbers and DOIs
should be taken from the publisher records when the reference list is
typeset; a few entries above carry venue and year but not pagination.
Nothing in the chapter's prose depends on those details.
